\section{The inequality and the proof}
\label{sec:rpi:introduction}

The Riemannian Penrose inequality compares the mass of an asymptotically
flat manifold with the total area of its outer minimizing minimal
boundary. Its geometric content is that a large minimal enclosure
requires a correspondingly large mass. Penrose's physical argument
\cite{Penrose1973} related an initial trapped area to the mass remaining
after gravitational collapse: cosmic censorship, area increase, and
settling to a Kerr black hole led to the proposed mass--area bound.
For time-symmetric initial data, the energy condition becomes
nonnegative scalar curvature and the horizon condition becomes
minimality. The resulting Riemannian problem retains the mass--area
comparison without requiring an evolution of the spacetime.

In three dimensions, one route from the horizon to infinity uses the Hawking mass
\cite{Hawking1968}. Geroch's calculation \cite{Geroch1973} showed its
monotonicity along smooth inverse mean curvature flow of two-spheres
in nonnegative scalar curvature. Jang and
Wald \cite[Section~5]{JangWald1977} derived the Penrose bound assuming
a global flow with the required asymptotics. Existence of such a flow
was the central difficulty: mean curvature may vanish and the smooth
evolution can develop singularities. Before the general proofs,
spherical symmetry yielded the inequality \cite{MalecOMurchadha1994},
and spinorial methods gave mass--area estimates involving an additional
normalized Sobolev quantity \cite{Herzlich1997}.

Huisken and Ilmanen's weak inverse mean curvature flow allows jumps
across the regions where smooth flow fails \cite{HuiskenIlmanen2001}.
It proves the sharp mass bound for each connected component of an
outermost minimal boundary in dimension three. Bray's conformal flow
\cite{Bray2001} instead changes the ambient metric while preserving the
total horizon area. Positive mass makes the evolving mass nonincreasing,
and the limiting comparison yields the sharp bound for the full area,
including disconnected boundaries. The harmonic-flat reduction,
area-preserving flow, and mass comparison in the present proof descend
from this construction.

Bray and Lee extended the conformal method to dimensions below eight
\cite{BrayLee2009}. The dimension restriction reflects the use of smooth
minimizing hypersurfaces. In
dimensions at least eight, a minimizing frontier may be singular, and
its singular points affect the flow construction, separation of slices,
and the geometric comparison used to decrease mass.

Restricted all-dimensional results were obtained earlier under
graphical hypotheses. Lam \cite[Corollary~7]{Lam2010Graphs} proved the
inequality for asymptotically flat Euclidean graphs with nonnegative
scalar curvature and convex horizon components. De Lima and Gir\~ao
\cite[Theorem~1.8]{deLimaGirao2015} treated two-sided asymptotically flat
quasi-graphs in Euclidean space, with nonnegative scalar curvature and
convex horizons in a horizontal hyperplane met orthogonally. These
results use additional graphical control of the geometry.

Bi and Zhu \cite[Theorem~1.2]{BiZhu2026v2} prove the all-dimensional
inequality together with Schwarzschild rigidity for almost-minimizing
Caccioppoli (finite-perimeter) boundaries whose regular part is minimal
and which are
their own outermost minimizing enclosures, with a $C^{2,\alpha}$ metric
extension through the boundary. Their version~2 conformal-flow argument
is the immediate framework for this paper. We develop the numerical
proof under stronger local minimization and ambient regularity
hypotheses: the frontier is locally perimeter minimizing on both sides,
the ambient metric is smooth through it, and the scalar curvature has
the pointwise decay stated below. Our initially outer minimizing
frontier is replaced by its outermost hull without changing its area.
The present contribution is a detailed numerical reconstruction
for this minimizing-frontier class.

The reconstruction makes explicit the preservation of area before
limiting nesting, uniform estimates near regular patches of tangent
cones, and complete smooth comparison metrics with positive conformal
factors. It uses only the largest limiting frontier radius to reach
the numerical conclusion. These reductions explain which geometric
and positive-mass inputs suffice without requiring convergence of every
frontier component or classification of equality.

\subsection{Manifolds, frontiers, and mass}

Write $n\ge3$, $k=n-1$, $d=n-2$, and
$\omega=\mathcal H^k_\delta(S^k)$. We use $O_j(r^{-\gamma})$ for a
tensor whose coordinate derivatives of order $\ell\le j$ are
$O(r^{-\gamma-\ell})$. On a Euclidean coordinate end the ADM mass is
\[
 m_{\rm ADM}(h)=\frac1{2k\omega}
 \lim_{R\to\infty}\int_{S_R}
 (\partial_jh_{ij}-\partial_ih_{jj})n_\delta^i\,dA_\delta.
\]
The decay and scalar integrability hypotheses below give the usual
well-defined mass in this normalization; see \cite{Bartnik1986}.

A possibly singular exterior is described by its filled side in a
smooth ambient manifold. The exterior contains the coordinate end; its
compact frontier $\Sigma$ is the support of the perimeter measure of
the filled side. Its area is the full reduced-boundary perimeter,
including every component. The metric extends smoothly through a
neighborhood of $\Sigma$. The exterior includes its closed frontier.
For $p,q\in E$, define the possibly extended length distance
\[
 d_E(p,q)=\inf\{\operatorname{Length}_h(\gamma):
       \gamma\subset E\text{ is a rectifiable path joining }p\text{ to }q\},
 \qquad \inf\varnothing=+\infty.
\]
Completeness means that each finite-distance equivalence class
$\{q\in E:d_E(p,q)<\infty\}$ is complete for the restricted metric.
This does not identify $E$ with the length completion of its open
exterior or assert finite-length access to every singular frontier
point. Every point and component of the full frontier remains included
in $E$, with its full perimeter counted.

The frontier is \emph{outer minimizing} if every finite-perimeter
enclosure of its filled side with compact additional support has at
least its perimeter. It is \emph{locally perimeter minimizing} if
compactly supported changes in sufficiently small neighborhoods,
including changes on either side, cannot reduce perimeter. The latter
condition is stronger than weak stationarity. It is the local
comparison needed in the flow construction. At free regular points
it gives the minimal-hypersurface equation.

\begin{theorem}[Riemannian Penrose inequality with minimizing frontiers]
\label{thm:rpi}
Let $(E,h)$ be a connected complete $n$-dimensional exterior, including
its compact frontier $\Sigma$ in the enclosing-set sense above. Suppose
that $h$ is smooth on $E$ and through $\Sigma$, and that
\begin{enumerate}
\item $R_h\ge0$, $R_h\in L^1(E,dV_h)$, and
      $R_h=O(r^{-n-\beta})$ for some $\beta>0$ on the end;
\item outside a compact set there is exactly one Euclidean coordinate
      end with $h-\delta=O_2(r^{-\tau})$, where $\tau>d/2$;
\item the full frontier is outer minimizing and locally perimeter
      minimizing as the boundary of the filled side. Its regular part
      is a smooth embedded minimal hypersurface, and its singular set
      has Hausdorff dimension at most $n-8$.
\end{enumerate}
Then
\begin{equation}\label{eq:rpi:numerical}
 m_{\rm ADM}(h)\ge\frac12
 \left(\frac{\mathcal H_h^k(\Sigma)}{\omega}\right)^{d/k}.
\end{equation}
\end{theorem}

All boundary components enter the displayed area. Neither spin nor
connectedness of the frontier is assumed. The zero-area case is
included: the density lower bound makes a nonempty locally minimizing
frontier have positive perimeter, so zero area means an empty
frontier and the conclusion is nonnegative mass.

For $3\le n\le7$, the frontier is smooth and the numerical assertion
follows from Bray--Lee's published Theorem~1.4 \cite{BrayLee2009}. Its
numerical part has no spin assumption; spin belongs to the separate equality
assertion there. The scalar decay required in that theorem holds here.
The rest of this paper develops the singular argument for $n\ge8$,
including the reduction from harmonic-flat ends to the stated weaker
asymptotics.

The constant is illustrated by the Schwarzschild--Tangherlini metric
\cite{Tangherlini1963}
\[
 h_m=\left(1+\frac{m}{2r^d}\right)^{4/d}\delta,
 \qquad r\ge(m/2)^{1/d},\quad m>0.
\]
Its inner sphere is minimal, its mass is $m$, and its area is
$\omega(2m)^{k/d}$. Thus it realizes the numerical constant in
Equation~\eqref{eq:rpi:numerical}.

The enclosure lemmas in Section~\ref{sec:g-enclosures} identify full
enclosing area with the area of a minimizing frontier and detach that
frontier from a strict inner barrier. Their geometric conclusions are
used independently in boundary-graph deformations
\cite{CompanionBoundaryGraphs}. Corollary~\ref{cor:g-strict-barrier-penrose}
combines them with Theorem~\ref{thm:rpi} to give a Penrose bound for smooth
strict inner barriers. This passage from enclosing area to a minimizing
frontier also supplies the numerical comparison in
\cite{CompanionSpacetimeNumerical} and the smooth nonnegative-mass
reduction in \cite{CompanionSpacetimeRigidity}.

\subsection{Four quantities that control the proof}

First assume the initial metric $g_0=h$ has a harmonic-flat end,
meaning $g_0=H^{4/d}\delta$ there with $H>0$ Euclidean harmonic and
$H\to1$. Extend $g_0$ smoothly a short distance behind its frontier
and truncate on that buried side at a smooth compact boundary. This
gives an ambient exterior $M$ and an initial filled region $K_0$
containing the buried boundary; the original end exterior is
$M\setminus K_0$. The extension needs no scalar-curvature sign behind
the original frontier.

The flow has metrics $g_t=u_t^{4/d}g_0$ and filled regions $K_t$ with
frontiers $\Sigma_t$ and end exteriors $E_t=M\setminus K_t$.
At infinity $u_t\to e^{-t}$. If $y$ denotes the original end coordinate,
we measure mass in the normalized coordinate $x=e^{-2t/d}y$,
in which the metric tends to $\delta$ at spatial infinity for each
fixed $t$.
Four quantities organize the proof:
\[
 A=P_{g_t}(K_t),\qquad m(t)=m_{\rm ADM}(g_t),\qquad
 \mu(t)=m(t)-c_t,\qquad R_t\quad\text{(outer radius)}.
\]
For the radius, represent the compact core by the ball of radius
$R_0e^{-2t/d}$ when the original chart starts at $|y|=R_0$; $R_t$ is
the maximum radius of the union of that ball and the frontier portion
in the normalized chart. The capacity
$c_t=\mathfrak c(\Sigma_t,g_t)$ in the deficit $\mu(t)$ is the
normalized energy of the harmonic function $w_t$
equal to one at the conductor and zero at infinity:
\[
 \mathfrak c(\Sigma_t,g_t)=\frac1{d\omega}
 \int_{E_t}|dw_t|_{g_t}^2\,dV_{g_t}.
\]
The flow keeps $A$ fixed and satisfies $m'=-2\mu$ almost everywhere.
The mass--capacity comparison gives $\mu\ge0$; a one-sided time estimate
then gives $\mu(t)\to0$. The limiting harmonic potential identifies
the largest limiting frontier radius, and enclosing
spheres give the sharp area inequality.

The constructions establishing these facts have a necessary order.
Section~\ref{sec:g-enclosures} states the independent enclosure results;
their proofs are collected in Section~\ref{sec:g-enclosure-proofs},
where they enter the final approximation argument. In
Sections~\ref{sec:rpi:flow-localization}--\ref{sec:rpi:flow-construction},
the discrete flow supplies uniform density and potential estimates.
We prove preservation of the limiting minimum area before deriving
nesting of the limiting outermost hulls. This follows the order in
Bi--Zhu version~2, Lemmas~2.21 and~2.23 and Corollary~2.22
\cite{BiZhu2026v2}.

For mass--capacity comparison we need a smooth boundary separated from
the singular frontier. Sections~\ref{sec:rpi:regular-patches}--\ref{sec:rpi:median-separation}
use a single median of the distance between frontiers to normalize all
regular patches and prove this separation. The method descends from the
singular maximum principles of Simon and Ilmanen
\cite{Simon1987Maximum,Ilmanen1996Maximum}, as developed in
\cite[Section~4]{Wickramasekera2014Maximum}, and adapts Bi--Zhu's relative
forcing and normalization \cite[Proposition~2.34 and Lemma~2.35]{BiZhu2026v2}.
Section~\ref{sec:rpi:smoothing} then smooths on the exterior side while
retaining the curvature sign, following the distance-offset strategy
of \cite{BiZhuSmoothing2026}.

Before applying positive mass, Sections~\ref{sec:rpi:smooth-repair} and
\ref{sec:rpi:lee-completion} construct complete smooth comparison metrics.
Brendle--Wang's theorem \cite[Corollary~1.6]{BrendleWang2026} requires
strictly positive scalar curvature and an all-order end expansion.
From it we obtain nonnegative mass on harmonic-flat ends and
mass--capacity comparison by doubling and collar repair \cite{Miao2002}.
A one-end reduction then permits the quantitative estimate of Lee
\cite{Lee2007NearEquality}. Appendix~\ref{sec:rpi:weighted-minimizers}
proves a local weighted-minimizer estimate using \cite{NaberValtorta2020}
to substantiate singular-set control in that positive-mass method.
The smooth positive-mass theorem itself remains an independent input.

Finally, Section~\ref{sec:rpi:numerical-limit} identifies the limiting
outer radius and completes the harmonic-flat argument.
Section~\ref{sec:rpi:density-repair} first creates a strict smooth inner
barrier and then approximates the general end, controlling both mass
and every enclosing area. Thus the approximation does not require a
theorem preserving a singular minimal frontier.

\paragraph{Mean-curvature convention.}
At an inner frontier, $\nu$ points toward the asymptotic end and
$H_\nu=\operatorname{div}_\Sigma\nu$. The outward normal of the exterior
region is $-\nu$. Mean-curvature signs below are always stated using
one of these explicit normals.
